% Capitolo autonomo di approfondimento.
% Non è incluso in main.tex, così non altera la struttura assegnata alla tesi.

\chapter{Approfondimento sullo Structural Similarity Index Measure}
\label{cap:approfondimento-ssim}

\section{Principio di funzionamento}

Lo Structural Similarity Index Measure (SSIM) è una metrica progettata per
confrontare due immagini tenendo conto di caratteristiche legate alla
percezione visiva. A differenza di una semplice differenza calcolata pixel per
pixel, SSIM considera tre componenti: luminanza, contrasto e struttura.

Date due porzioni di immagine \(x\) e \(y\), una forma compatta della metrica è:

\begin{equation}
    \operatorname{SSIM}(x,y)=
    \frac{(2\mu_x\mu_y+C_1)(2\sigma_{xy}+C_2)}
    {(\mu_x^2+\mu_y^2+C_1)(\sigma_x^2+\sigma_y^2+C_2)} ,
    \label{eq:ssim}
\end{equation}

dove \(\mu_x\) e \(\mu_y\) rappresentano le medie dei valori dei pixel,
\(\sigma_x^2\) e \(\sigma_y^2\) le rispettive varianze e \(\sigma_{xy}\) la
covarianza tra le due porzioni. Le costanti \(C_1\) e \(C_2\) stabilizzano il
calcolo quando i denominatori assumono valori molto piccoli.

La metrica viene normalmente calcolata su finestre locali e i risultati sono
poi aggregati per ottenere uno score complessivo. Il valore massimo è \(1\) e
indica immagini identiche rispetto al confronto effettuato. Sebbene nei casi
più comuni lo score osservato sia compreso tra \(0\) e \(1\), il suo intervallo
teorico può includere anche valori negativi.

\section{Confronto con le metriche pixel per pixel}

Metriche come il Mean Squared Error misurano direttamente la distanza tra i
valori dei pixel corrispondenti. Una piccola variazione di illuminazione, una
sfocatura o uno spostamento della fotocamera possono quindi produrre una
differenza numerica elevata anche quando il contenuto della scena è ancora
riconoscibile.

SSIM valuta invece la conservazione della struttura locale. Nell'esempio
seguente, la seconda immagine contiene lo stesso soggetto della prima ma è
sfocata: i valori dei singoli pixel cambiano, mentre gran parte della struttura
visiva resta simile.

\begin{center}
    \begin{minipage}{0.75\linewidth}
        \centering
        \IfFileExists{Risorse/SSIM1.png}
            {\includegraphics[width=\linewidth]{Risorse/SSIM1.png}}
            {\fbox{\parbox[c][4cm][c]{0.9\linewidth}{
                \centering Inserire qui il confronto tra l'immagine originale
                e la versione sfocata.
            }}}
        \captionof{figure}{Confronto tra un'immagine originale e una sua
        versione sfocata.}
        \label{fig:ssim-originale-sfocata}
    \end{minipage}
\end{center}

Lo score non stabilisce quale delle due immagini sia migliore e non riconosce
gli oggetti rappresentati. Esprime soltanto quanto le loro strutture visive
siano simili secondo la metrica.

\begin{center}
    \begin{minipage}{0.75\linewidth}
        \centering
        \IfFileExists{Risorse/SSIM2.png}
            {\includegraphics[width=\linewidth]{Risorse/SSIM2.png}}
            {\fbox{\parbox[c][4cm][c]{0.9\linewidth}{
                \centering Inserire qui l'esempio degli score ottenuti dal
                confronto.
            }}}
        \captionof{figure}{Esempio di interpretazione dello score SSIM.}
        \label{fig:ssim-confronto-score}
    \end{minipage}
\end{center}

\section{Impiego nel progetto}

Nella modalità automatica di VisionCaption, SSIM viene usato come primo livello
di selezione dei frame. Il sistema confronta il frame corrente con un frame di
riferimento per stimare se la scena sia cambiata in misura sufficiente da
giustificare una nuova elaborazione.

Se la somiglianza rimane sopra la soglia configurata, il frame può essere
scartato senza richiedere nuovamente le fasi più costose della pipeline. Se la
somiglianza scende sotto la soglia, il sistema prosegue con i controlli
successivi, tra cui il rilevamento degli oggetti. In questo modo si evita, per
esempio, di generare ripetutamente la stessa descrizione mentre l'utente
inquadra per diversi secondi una scena sostanzialmente immobile.

Il costo del confronto cresce linearmente rispetto al numero di pixel
analizzati, cioè è approssimativamente \(\mathcal{O}(n)\). Questa caratteristica
lo rende adatto a fungere da filtro leggero prima dei modelli di visione più
onerosi.

\section{Limiti e integrazione con il rilevamento semantico}

SSIM individua una variazione visiva, ma non ne comprende il significato. Un
movimento della fotocamera, un cambiamento di luminosità, una sfocatura o
un'occlusione temporanea possono ridurre lo score senza che nella scena sia
comparso un elemento realmente importante. Al contrario, una variazione
semanticamente rilevante ma circoscritta a una piccola area potrebbe produrre
solo una modifica limitata dello score complessivo.

La soglia deve quindi essere scelta sperimentalmente in relazione alla
risoluzione dei frame e alle condizioni d'uso; non esiste un valore universale
adatto a ogni situazione. Anche la scelta del frame di riferimento è
significativa: il confronto con il frame precedente è sensibile ai cambiamenti
graduali, mentre il confronto con un frame chiave può accumulare nel tempo
piccole variazioni.

Per questi motivi, nel progetto SSIM non è impiegato come unico criterio
decisionale. La strategia ibrida lo combina con il rilevamento degli oggetti:
SSIM fornisce un'indicazione rapida del cambiamento visivo, mentre il rilevatore
aggiunge informazioni sul contenuto della scena. La combinazione riduce sia le
caption ridondanti sia le attivazioni causate da variazioni prive di interesse
semantico.
